\section*{Hoofdstuk \ref{ch:g4:symmetry} --- Symmetrie}

\begin{solution}{exo:g4:symmetry:1}
Vierkant: ja ($4$ assen). Rechthoek: ja ($2$). Ongelijkzijdige driehoek:
geen enkele vouw werkt. Cirkel: ja --- elke lijn door het middelpunt.
\end{solution}

\begin{solution}{exo:g4:symmetry:2}
Vierkant: $4$ (twee middellijnen, twee diagonalen). Rechthoek: $2$ (alleen
de middellijnen). Gelijkzijdige driehoek: $3$ (één door elk \omterm{def:g4:geometry:polygon}{hoekpunt}).
\end{solution}

\begin{solution}{exo:g4:symmetry:3}
Vouw je langs de diagonaal, dan hebben de twee driehoekige \omterm{def:g3:division:half}{helften} dezelfde
vorm, maar liggen ze niet op elkaar --- de ene ligt ``verkeerd om'': de
overtrek komt naast de figuur terecht en niet erop. De diagonaal is dus geen
as.
\end{solution}

\begin{solution}{exo:g4:symmetry:4}
De beeld-L staat $2$ vakjes aan de \emph{andere} kant van de as, in
spiegelschrift (een omgekeerde L), met dezelfde maten.
\end{solution}

\begin{solution}{exo:g4:symmetry:5}
De afgemaakte figuur toont het vlaggetje met daaronder zijn spiegelbeeld
onder de as, mast onder mast.
\end{solution}

\begin{solution}{exo:g4:symmetry:6}
$P$ beweegt niet: bij het vouwen langs de as blijft elk punt \emph{op} de as
precies waar het is --- $P$ is zijn eigen beeld.
\end{solution}

\begin{solution}{exo:g4:symmetry:7}
Verticale as: A, H, O, T. Horizontale as: B, E, H, O. Beide: H, O. Geen van
beide: F, N.
\end{solution}

\begin{solution}{exo:g4:symmetry:8}
Ja: de driehoek is gelijkbenig, en zijn as loopt door de top tussen de twee
zijden van $5$ cm en door het midden van de basis van $3$ cm. De vouw wisselt
de twee gelijke zijden om.
\end{solution}

\begin{solution}{exo:g4:symmetry:9}
Elke figuur met precies twee assen is goed: een rechthoek die geen vierkant
is, een plusteken met armen van twee verschillende lengtes (lang horizontaal,
kort verticaal), of een ovaal. De twee assen snijden elkaar altijd in het
midden van de figuur.
\end{solution}

\begin{solution}{exo:g4:symmetry:10}
Het hele huis is de linkerhelft plus haar spiegeltweeling: een muur die twee
keer zo breed is als de halve muur, onder een volledig driehoekig dak. Elke
lengte rechts is gelijk aan haar partner links --- spiegelen kopieert lengtes
precies.
\end{solution}

\begin{solution}{exo:g4:symmetry:11}
De tekeningen steunen beide richtingen: de gelijkbenige driehoek vouwt langs
de lijn door zijn top (twee gelijke zijden $\to$ een as); de gelijkzijdige
vouwt op drie manieren; de ongelijkzijdige heeft geen as (geen gelijke zijden
$\to$ geen as). Aya's dubbele bewering is juist --- ze wordt bewezen in
\cref{ch:g6:symmetry}.
\end{solution}
